@ID: ON20D1P1U
@BIDANG: Kombinatorika

Bila $x, y$ dan $m$ adalah bilangan bulat dengan $x,y\ge0$ dan $m\ge x+y$, buktikan bahwa

$$
\binom{m+1}{x+y+1}
=
\sum_{i=0}^{m}
\binom{i}{x}
\binom{m-i}{y}.
$$

@ID: ON20D1P2U
@BIDANG: Analisis Kompleks

Tentukan semua fungsi $f:\mathbb C\rightarrow\mathbb C$ yang memenuhi $f(0)=0$ dan

$$
|f(z)-f(w)|=|z-w|
$$

untuk setiap $z\in\mathbb C$ dan $w\in\{0,1,i\}$.

@ID: ON20D1P3U
@BIDANG: Analisis Real

Diketahui fungsi $f:[a,b]\rightarrow\mathbb R$ kontinu dan

$$
F(x)=\int_a^x f(t)\,dt,
$$

untuk setiap $x\in[a,b]$. Jika $(p_n)\subseteq\mathbb R$ merupakan barisan Cauchy dengan $a\le p_n\le b$ untuk setiap $n$, buktikan bahwa terdapat $p$ dan $c$, dengan $a\le p,c\le b$, sehingga barisan $(F(p_n))$ konvergen ke

$$
f(c)(p-a).
$$

@ID: ON20D1P4U
@BIDANG: Aljabar Linear

Misalkan $V$ sebuah ruang vektor berdimensi $n\ge2$ dan $\mathcal B=\{v_1,v_2,\dots,v_n\}$ sebuah basis dari $V$. Diberikan sebuah operator linear invertibel $T:V\rightarrow V$ dengan balikan $T^{-1}$.

Operator linear $A$ dan $B$ pada $V$ didefinisikan oleh

- $A(v_1)=0$ dan $A(v_i)=T^{-1}(v_i)$ untuk $i=2,\dots,n$.
- $B(v_1)=0$ dan $B(v_i)=T^{-1}(v_{i-1})$ untuk $i=2,\dots,n$.

(a) Buktikan bahwa

$$
\operatorname{rank}\big((A\circ T)^{\,n-1}\big)=n-1.
$$

(b) Hitunglah

$$
\operatorname{rank}\big((B\circ T)^{\,n-1}\big).
$$

Catatan: Jika $S$ suatu operator linear pada $V$ dan $k$ suatu bilangan asli, maka

$$
S^k=\underbrace{S\circ S\circ\cdots\circ S}_{k\text{ kali}}.
$$

@ID: ON20D1P5U
@BIDANG: Struktur Aljabar

Misalkan $G$ adalah suatu grup hingga dan $m$ adalah suatu bilangan asli kelipatan tiga. Buktikan jika $G$ memuat tepat sebanyak $m$ unsur berorde $m$, maka $m$ adalah bilangan kelipatan enam dan $G$ memiliki tepat tiga subgrup siklik berorde $m$.

@ID: ON20D2P1U
@BIDANG: Analisis Kompleks

Pasangan bilangan kompleks $(a,b)$ dikatakan seimbang jika

$$
|a|=|b|=|a+b|.
$$

(a) Tunjukkan bahwa terdapat pasangan seimbang $(a,b)$ dengan $a\ne b$.

(b) Jika $(a,b)$ merupakan pasangan seimbang dan $a^{2020}=b^{2020}$, tunjukkan bahwa $a=b$.

@ID: ON20D2P2U
@BIDANG: Kombinatorika

Andaikan $B=\{1,2,\dots,n\}$. Untuk bilangan bulat positif $k\le n$, buktikan bahwa banyaknya barisan

$$
B_1,B_2,\dots,B_k
$$

sedemikian sehingga $B_i\subseteq B$ untuk $1\le i\le k$ dan

$$
\bigcup_{i=1}^k B_i=B
$$

adalah

$$
(2^k-1)^n.
$$

@ID: ON20D2P3U
@BIDANG: Analisis Real

Diberikan $a\in\mathbb R$ dan interval terbuka terbatas $I\subseteq\mathbb R$ dengan $a\in I$. Untuk setiap $n\in\mathbb N$, fungsi $f_n,g,h:I\rightarrow\mathbb R$ terdiferensial dan memenuhi:

(a) $f_n(a)=g(a)$ untuk setiap $n\in\mathbb N$;

(b) $f_n'(x)\ge g'(x)$ untuk setiap $x\in I$ dengan $x>a$;

(c) $f_n'\to h'$ secara seragam pada $I$.

Buktikan bahwa terdapat fungsi terdiferensial $f:I\rightarrow\mathbb R$ sehingga $f_n\to f$ pada $I$ dan

$$
h(x)\ge g(x)
$$

untuk setiap $x\in I$ dengan $x>a$.

@ID: ON20D2P4U
@BIDANG: Struktur Aljabar

Misalkan $R$ adalah ring yang memiliki identitas perkalian $1$. Diketahui bahwa $x,y\in R$ memenuhi

$$
xy=1
\quad\text{tetapi}\quad
yx\ne1.
$$

(a) Tunjukkan bahwa untuk setiap bilangan bulat positif $m\ne n$ berlaku

$$
y^mx^m\ne y^nx^n.
$$

(b) Tunjukkan bahwa terdapat tak berhingga banyaknya unsur $z\in R$ yang memenuhi

$$
z^{2020}=0.
$$

@ID: ON20D2P5U
@BIDANG: Aljabar Linear

Misalkan $A,B,C,$ dan $D$ adalah matriks berukuran $9\times9$ dengan entri-entri bilangan real dan memenuhi

$$
CA=BC.
$$

Jika polinom karakteristik matriks $A$ sama dengan polinom karakteristik matriks $B$, apakah polinom karakteristik matriks $(A-DC)$ juga sama dengan polinom karakteristik matriks $(B-CD)$?

Lengkapilah jawaban Anda dengan bukti.